%! Linear Transformations and Matrices — Unit 8, Lesson 8.8 — Lesson Plan
\documentclass[10pt]{article}
\usepackage{precalculus-article}
\usepackage{precalculus-boxes}
\usepackage{graphicx}
\graphicspath{{images/}}
\newcommand{\TallMath}[1]{$\displaystyle #1\rule[-1.4em]{0pt}{3.2em}$}
\newcommand{\CourseName}{Precalculus}
\newcommand{\MeetingLength}{60 minutes}
\newcommand{\UnitNumberName}{Unit 8: Parametric Functions, Vectors, and Matrices \quad}
\newcommand{\LessonNumberName}{Lesson 8.8: Linear Transformations and Matrices}

\begin{document}
\begin{center}
    {\Large\bfseries \CourseName \\
    {\small\bfseries \UnitNumberName \LessonNumberName}}
\end{center}

\begin{tcolorbox}[colback=sky, colframe=navy, boxrule=0.9pt, arc=2mm,
    left=2mm, right=2mm, top=1.5mm, bottom=1.5mm]
    \textbf{Primary Objective:} Students will recognize linear transformations as
    matrix-vector multiplications, compute output vectors using $L(\vec{v}) = A\vec{v}$,
    and apply a $2\times 2$ transformation matrix to multiple input vectors simultaneously.
    (Big Idea: Functions)
\end{tcolorbox}

\renewcommand{\arraystretch}{1.6}
\begin{skillbox}[Priority Ideas \& Skills]{goldbox}
\begin{minipage}[t]{0.46\linewidth}\vspace{0pt}
    \textbf{AP Skill:} 1.B --- Express equivalent analytic forms\\[4pt]
    \textbf{Essential Knowledge:}\\
    EK 4.12.A.1 -- 4.12.A.5
\end{minipage}
\hfill
\begin{minipage}[t]{0.50\linewidth}\vspace{0pt}
    \textbf{Key Understandings:}
    \begin{itemize}[leftmargin=1.2em, itemsep=1pt]
        \item A linear transformation maps each output component as a sum of constant
              multiples of the input components.
        \item Any linear transformation $L:\mathbb{R}^2\to\mathbb{R}^2$ is represented
              by a unique $2\times 2$ matrix $A$ via $L(\vec{v})=A\vec{v}$.
        \item Applying $A$ to a $2\times n$ matrix transforms $n$ vectors simultaneously.
    \end{itemize}
\end{minipage}
\end{skillbox}

\begin{skillbox}[{Vocabulary, Concepts, \& Theorems}]{greenbox}
\begin{itemize}[leftmargin=1.2em, itemsep=2pt]
    \item \textbf{Linear transformation} --- a function $L:\mathbb{R}^2\to\mathbb{R}^2$
          where each output component is a sum of constant multiples of input components.
    \item \textbf{Transformation matrix} --- the unique $2\times 2$ matrix $A$ such that
          $L(\vec{v})=A\vec{v}$ for all $\vec{v}\in\mathbb{R}^2$.
    \item \textbf{Input vector / output vector} --- $\vec{v}$ (domain) and $L(\vec{v})=A\vec{v}$
          (range).
\end{itemize}
\end{skillbox}

\begin{skillbox}[Activate Prior Knowledge \& Spiral Review]{sky}
    \begin{tabularx}{\linewidth}{@{}X c@{}}
        \begin{minipage}[t]{0.55\linewidth}\vspace{0pt}
            Please see attached warm-up. Prior skills reviewed:
            \begin{itemize}[leftmargin=1.2em, itemsep=2pt]
                \item Matrix multiplication (Lesson 4.10)
                \item Vectors as $2\times 1$ column matrices (Lesson 4.9)
                \item The zero vector
            \end{itemize}
        \end{minipage}
        &
        \begin{minipage}[t]{0.38\linewidth}\vspace{0pt}\centering
        \end{minipage}
    \end{tabularx}
\end{skillbox}

\begin{skillbox}[Hook]{sky}
    Computer graphics moves points on screen by transforming their position vectors.
    If a triangle has vertices at $(1,0)$, $(0,1)$, and $(-1,0)$, can we scale all
    three vertices at once using a single matrix multiplication?  Today we discover the
    connection between matrix multiplication and geometric transformations.
\end{skillbox}

\begin{skillbox}[Lesson]{sky}
\begin{multicols}{2}

\textbf{Step 1 --- Definition of a Linear Transformation.}\\
A function $L:\mathbb{R}^2\to\mathbb{R}^2$ is a \emph{linear transformation} if each
output component is a sum of constant multiples of the input components:
\[
  L\!\left(\begin{bmatrix}x\\y\end{bmatrix}\right)
  = \begin{bmatrix}ax+by\\cx+dy\end{bmatrix}.
\]
Key consequence: $L(\vec{0}) = \vec{0}$ (the zero vector maps to the zero vector).

\vspace{6pt}
\textbf{Step 2 --- Matrix Representation.}\\
The constants $a,b,c,d$ can be organized into a $2\times 2$ matrix
$A = \begin{bmatrix}a & b \\ c & d\end{bmatrix}$, so that
\[
  L(\vec{v}) = A\vec{v}.
\]
For \emph{any} $2\times 2$ matrix $A$, the function $L(\vec{v})=A\vec{v}$ is a linear
transformation.  Conversely, every $L$ has exactly one such matrix.

\vspace{6pt}
\textbf{Step 3 --- Zero Vector Check.}\\
Since matrix multiplication is distributive, $A\vec{0} = \vec{0}$.  A mapping that does
\emph{not} send $\vec{0}\to\vec{0}$ cannot be linear.

\columnbreak

\textbf{Step 4 --- Batch Transformation.}\\
Expressing $n$ vectors as columns of a $2\times n$ matrix and multiplying by $A$ gives
all $n$ output vectors at once:
\[
  A\begin{bmatrix}\vec{v}_1 & \vec{v}_2 & \cdots & \vec{v}_n\end{bmatrix}
  = \begin{bmatrix}A\vec{v}_1 & A\vec{v}_2 & \cdots & A\vec{v}_n\end{bmatrix}.
\]

\vspace{6pt}
\textbf{Worked Example --- Scaling Transformation.}\\
Let $A = \begin{bmatrix}2&0\\0&2\end{bmatrix}$.  Transform the triangle vertices
$(1,0)$, $(0,1)$, $(-1,0)$ simultaneously:
\[
  \begin{bmatrix}2&0\\0&2\end{bmatrix}
  \begin{bmatrix}1&0&-1\\0&1&0\end{bmatrix}
  = \begin{bmatrix}2&0&-2\\0&2&0\end{bmatrix}.
\]
Each vertex is scaled by a factor of 2.  The transformation is a dilation from the origin.

\end{multicols}
\end{skillbox}

\begin{skillbox}[Active Monitoring]{redbox}
    \textbf{Circulate and check:}
    \begin{itemize}[leftmargin=1.2em, itemsep=2pt]
        \item Are students multiplying row $\times$ column correctly (not column $\times$ column)?
        \item Do students recognize that $L(\vec{0})=\vec{0}$ is a \emph{necessary} condition,
              not sufficient?
        \item Cold-call: ``What is the transformation matrix for the identity transformation?''
        \item Cold-call: ``If I told you $L$ maps $(1,0)\to(3,5)$ and $(0,1)\to(-1,2)$,
              what is $A$?''
    \end{itemize}
\end{skillbox}

\begin{skillbox}[Group Work \& Differentiation]{redbox}
    \textbf{Tier R (Remediate):} Given $A=\begin{bmatrix}2&0\\0&2\end{bmatrix}$, compute
    $L(\vec{v})=A\vec{v}$ for $\vec{v}=\langle 3,1\rangle$ and $\vec{v}=\langle -1,4\rangle$.
    Describe the transformation in words.\\[4pt]
    \textbf{Tier A (At Standard):} Given $A=\begin{bmatrix}0&-1\\1&0\end{bmatrix}$, apply $L$
    to the triangle vertices $(1,0)$, $(0,1)$, $(-1,-1)$. Describe the geometric effect.\\[4pt]
    \textbf{Tier E (Extend):} Show that $L(\vec{v})=A\vec{v}$ with $A=\begin{bmatrix}1&2\\0&1\end{bmatrix}$
    is linear by verifying $L(\vec{v}+\vec{w})=L(\vec{v})+L(\vec{w})$ for two specific vectors.
\end{skillbox}

\begin{skillbox}[Individual Work \& Assessment]{redbox}
    \textbf{Exit Ticket (last 5 minutes):}
    \begin{enumerate}[leftmargin=1.5em, itemsep=2pt]
        \item Let $A=\begin{bmatrix}3&1\\0&2\end{bmatrix}$.  Find $L\!\left(\begin{bmatrix}2\\1\end{bmatrix}\right)$.
        \item A linear transformation maps $\begin{bmatrix}1\\0\end{bmatrix}\to\begin{bmatrix}3\\5\end{bmatrix}$
              and $\begin{bmatrix}0\\1\end{bmatrix}\to\begin{bmatrix}-1\\2\end{bmatrix}$.  Write matrix $A$.
    \end{enumerate}
    \textbf{Assessment note:} Q2 assesses whether students understand that the columns of $A$ are the
    images of the standard basis vectors --- a key insight from EK 4.12.A.4.
\end{skillbox}

\begin{skillbox}[Reinforcement \& Extension]{goldbox}
    \textbf{Homework:} 5 problems covering single-vector transformation, batch transformation,
    identifying a transformation geometrically, linearity check (zero-vector test), and an
    AP-style MC item.\\[4pt]
    \textbf{Preview of Lesson 4.13:} Matrices as functions --- composition of transformations
    and the connection to matrix multiplication of two transformation matrices.
\end{skillbox}
\end{document}
